\documentclass[otchet]{SCWorks}
% Тип обучения (одно из значений):
%    bachelor   - бакалавриат (по умолчанию)
%    spec       - специальность
%    master     - магистратура
% Форма обучения (одно из значений):
%    och        - очное (по умолчанию)
%    zaoch      - заочное
% Тип работы (одно из значений):
%    coursework - курсовая работа (по умолчанию)
%    referat    - реферат
%  * otchet     - универсальный отчет
%  * nirjournal - журнал НИР
%  * digital    - итоговая работа для цифровой кафедры
%    diploma    - дипломная работа
%    pract      - отчет о научно-исследовательской работе
%    autoref    - автореферат выпускной работы
%    assignment - задание на выпускную квалификационную работу
%    review     - отзыв руководителя
%    critique   - рецензия на выпускную работу

% * Добавлены вручную. За вопросами к @mchernigin
\usepackage{preamble}

\begin{document}

% Кафедра (в родительном падеже)
\chair{математической кибернетики и компьютерных наук}

% Тема работы
\title{Анализ сложности для сортировок, не использующих сравнение элементов}

% Курс
\course{2}

% Группа
\group{211}

% Факультет (в родительном падеже) (по умолчанию "факультета КНиИТ")
% \department{факультета КНиИТ}

% Специальность/направление код - наименование
 \napravlenie{02.03.02 "--- Фундаментальная информатика и информационные технологии}
% \napravlenie{02.03.01 "--- Математическое обеспечение и администрирование информационных систем}
% \napravlenie{09.03.01 "--- Информатика и вычислительная техника}
%\napravlenie{09.03.04 "--- Программная инженерия}
% \napravlenie{10.05.01 "--- Компьютерная безопасность}

% Для студентки. Для работы студента следующая команда не нужна.
% \studenttitle{студентки}

% Фамилия, имя, отчество в родительном падеже
\author{Бурова Арсения Александровича}

% Заведующий кафедрой 
\chtitle{доцент, к.\,ф.-м.\,н.}
\chname{С.\,В.\,Миронов}

% Руководитель ДПП ПП для цифровой кафедры (перекрывает заведующего кафедры)
% \chpretitle{
%     заведующий кафедрой математических основ информатики и олимпиадного\\
%     программирования на базе МАОУ <<Ф"=Т лицей №1>>
% }
% \chtitle{г. Саратов, к.\,ф.-м.\,н., доцент}
% \chname{Кондратова\, Ю.\,Н.}

% Научный руководитель (для реферата преподаватель проверяющий работу)
\satitle{доцент, к.\,ф.-м.\,н.} %должность, степень, звание
\saname{Ю.\,Н.\,Кондратова}

% Руководитель практики от организации (руководитель для цифровой кафедры)
\patitle{доцент, к.\,ф.-м.\,н.}
\paname{С.\,В.\,Миронов}

% Руководитель НИР
\nirtitle{доцент, к.\,п.\,н.} % степень, звание
\nirname{В.\,А.\,Векслер}

% Семестр (только для практики, для остальных типов работ не используется)
\term{2}

% Наименование практики (только для практики, для остальных типов работ не
% используется)
\practtype{учебная}

% Продолжительность практики (количество недель) (только для практики, для
% остальных типов работ не используется)
\duration{2}

% Даты начала и окончания практики (только для практики, для остальных типов
% работ не используется)
\practStart{01.07.2024}
\practFinish{13.01.2024}

% Год выполнения отчета
\date{2026}

\maketitle

% Включение нумерации рисунков, формул и таблиц по разделам (по умолчанию -
% нумерация сквозная) (допускается оба вида нумерации)
\secNumbering

\tableofcontents

% Раздел "Обозначения и сокращения". Может отсутствовать в работе
% \abbreviations
% \begin{description}
%     \item ... "--- ...
%     \item ... "--- ...
% \end{description}

% Раздел "Определения". Может отсутствовать в работе
% \definitions

% Раздел "Определения, обозначения и сокращения". Может отсутствовать в работе.
% Если присутствует, то заменяет собой разделы "Обозначения и сокращения" и
% "Определения"
% \defabbr

\section{Реализация сортировок}
\subsection{Сортировка подсчетом}
\begin{Verbatim}[
  numbers=left,
  breaklines=true,
  breakanywhere=true
]
void countingSort(int A[], int n, int k) {
    //массив счетчиков размера k, заполняем 0
    int* C = new int[k]();
    //количество каждого числа
    for (int i = 0; i < n; i++) {
        C[A[i]]++;
    }
    //восстанавливаем отсортированный массив
    int pos = 0;
    for (int number = 0; number < k; number++) {
        for (int i = 0; i < C[number]; i++) {
            A[pos] = number;
            pos++;
        }
    }
    delete[] C;
}
\end{Verbatim}

\subsection{Поразрядная сортировка LSD}

Функция radixSortLSD:
\begin{Verbatim}[
  numbers=left,
  breaklines=true,
  breakanywhere=true
]
void radixSortLSD(int A[], int N, int m, int P = 10) {
    //создаем P дополнительных массивов
    int** dop = new int* [P];
    for (int i = 0; i < P; i++) {
        dop[i] = new int[N];
    }

    //массив размеров дополнительных массивов
    int* dopSizes = new int[P]();

    //проходим по всем разрядам от 1 до m
    for (int i = 1; i <= m; i++) {
        //очищаем доп массивы
        for (int b = 0; b < P; b++) {
            dopSizes[b] = 0;
        }

        //распределяем элементы по доп массивам
        for (int j = 0; j < N; j++) {
            int d = digit(A[j], i);
            dop[d][dopSizes[d]] = A[j];
            dopSizes[d]++;
        }

        //перезаписываем массив A из доп массивов
        int pos = 0;
        for (int b = 0; b < P; b++) {
            for (int j = 0; j < dopSizes[b]; j++) {
                A[pos] = dop[b][j];
                pos++;
            }
        }
    }

    for (int i = 0; i < P; i++) {
        delete[] dop[i];
    }
    delete[] dop;
    delete[] dopSizes;
}
\end{Verbatim}

Функция findMaxRazryad:
\begin{Verbatim}[
  numbers=left,
  breaklines=true,
  breakanywhere=true
]
int findMaxRazryad(int arr[], int n) {
    int maxVal = arr[0];
    for (int i = 1; i < n; i++) {
        if (arr[i] > maxVal) maxVal = arr[i];
    }
    int razryad = 0;
    while (maxVal > 0) {
        maxVal /= 10;
        razryad++;
    }
    return razryad;
}
\end{Verbatim}

Функция digit:
\begin{Verbatim}[
  numbers=left,
  breaklines=true,
  breakanywhere=true
]
int digit(int x, int i) {
    for (int p = 1; p < i; p++) {
        x /= 10;
    }
    return x % 10;
}
\end{Verbatim}


\newpage

\section{Анализ сложности алгоритмов}

Перед вызовом функции \texttt{countingSort} происходит вычисление параметра $k$ с помощью функции \texttt{findk}. Функция \texttt{findk} находит максимальное значение в массиве за время $O(n)$, где $n$ — размер исходного массива, так как выполняется один проход по всем элементам для поиска максимума.

Сама функция \texttt{countingSort(int A[], int n, int k)} получает уже готовый параметр $k = \texttt{maxVal} + 1$.

Создание массива счетчиков \texttt{int* C = new int[k]()} и его инициализация нулями выполняется за время $O(k)$, где $k$ — размер массива счетчиков (равный максимальному значению элемента массива плюс единица).

Первый цикл для подсчёта частот элементов в массиве \texttt{C}, выполняется за время $O(n)$, так как каждый элемент массива \texttt{A} обрабатывается один раз:
\begin{verbatim}
for (int i = 0; i < n; i++) {
    C[A[i]]++;
}
\end{verbatim}


Восстановление отсортированного массива выполняется во втором цикле: внешний цикл перебирает все возможные значения от $0$ до $k-1$, а вложенный цикл записывает каждое число столько раз, сколько оно встречалось в исходном массиве.
\begin{verbatim}
for (int number = 0; number < k; number++) {
    for (int i = 0; i < C[number]; i++) {
        A[pos] = number;
        pos++;
    }
}
\end{verbatim}
Внешний цикл выполняется $k$ итераций. Внутренний цикл суммарно по всем итерациям внешнего цикла выполняет $n$ операций записи, так как каждый элемент исходного массива записывается ровно один раз. Таким образом, восстановление массива работает за время $O(n + k)$.


\textbf{Итог:}
\begin{itemize}
    \item $n$ — количество элементов в исходном массиве.
    \item $k$ — параметр, равный $\texttt{maxVal} + 1$ (максимальное значение элемента массива плюс единица).
    \item Функция \texttt{findk}: $O(n)$.
    \item Функция \texttt{countingSort}: $O(n + k)$.
    \item \textbf{Общая сложность:} $O(n + k)$.
\end{itemize}

\subsection{Анализ сложности сортировки подсчетом}
Функция \texttt{findMaxRazryad(arr, n)} для определения максимального количества разрядов в числах массива выполняется за время $O(n + d)$, где $n$ — размер исходного массива, а $d$ — максимальное количество разрядов в числе. Сначала происходит поиск максимального элемента за $O(n)$, затем подсчёт его разрядов за $O(d)$.

Создание массива корзин (10 дополнительных массивов) и массива для хранения их размеров происходит за время $O(1)$, так как количество корзин является константой и не зависит от размера входных данных.

Внешний цикл проходит по разрядам от младшего к старшему и выполняется $d$ раз, где $d$ --- число разрядов в числах.

Очистка массивов размеров корзин перед распределением выполняется за время $O(P) = O(1)$, так как $P = 10$ — константа.

Во внутреннем цикле каждое число помещается в соответствующую корзину, что осуществляется за $O(n)$ на каждой итерации. Для определения текущей цифры числа вызывается функция \texttt{digit}, работающая за $O(1)$ (количество разрядов ограничено константой).

Сборка чисел из корзин обратно в исходный массив осуществляется за $O(n)$ на каждой итерации, так как суммарно через все корзины проходит $n$ элементов.

Таким образом, каждый проход по одному разряду выполняется за $O(n)$. Всего за работу алгоритма получим $d$ проходов, следовательно, общая временная сложность составляет:

\[
O(n + d) + O(1) + O(n \cdot d) + O(n \cdot d) = O(n \cdot d)
\]

\textbf{Итог:}
\begin{itemize}
    \item $n$ — количество элементов в массиве.
    \item $d$ — максимальное количество разрядов в числах массива
    \item Функция \texttt{findMaxRazryad}: $O(n + d)$.
    \item Функция \texttt{radixSortLSD}: $O(n \cdot d)$.
    \item \textbf{Общая временная сложность:} $O(n \cdot d)$.
\end{itemize}
\newpage

\section{Приложение 1. Код программы для задания}

\begin{Verbatim}[
  numbers=left,
  breaklines=true,
  breakanywhere=true
]
#include <iostream>

using namespace std;

void countingSort(int A[], int n, int k) {
    //массив счетчиков размера k, заполняем 0
    int* C = new int[k]();
    //количество каждого числа
    for (int i = 0; i < n; i++) {
        C[A[i]]++;
    }
    //восстанавливаем отсортированный массив
    int pos = 0;
    for (int number = 0; number < k; number++) {
        for (int i = 0; i < C[number]; i++) {
            A[pos] = number;
            pos++;
        }
    }
    delete[] C;
}

int digit(int x, int i) {
    for (int p = 1; p < i; p++) {
        x /= 10;
    }
    return x % 10;
}


void radixSortLSD(int A[], int N, int m, int P = 10) {
    //создаем P дополнительных массивов
    int** dop = new int* [P];
    for (int i = 0; i < P; i++) {
        dop[i] = new int[N];
    }

    //массив размеров дополнительных массивов
    int* dopSizes = new int[P]();

    //проходим по всем разрядам от 1 до m
    for (int i = 1; i <= m; i++) {
        //очищаем доп массивы
        for (int b = 0; b < P; b++) {
            dopSizes[b] = 0;
        }

        //распределяем элементы по доп массивам
        for (int j = 0; j < N; j++) {
            int d = digit(A[j], i);
            dop[d][dopSizes[d]] = A[j];
            dopSizes[d]++;
        }

        //перезаписываем массив A из доп массивов
        int pos = 0;
        for (int b = 0; b < P; b++) {
            for (int j = 0; j < dopSizes[b]; j++) {
                A[pos] = dop[b][j];
                pos++;
            }
        }
    }

    for (int i = 0; i < P; i++) {
        delete[] dop[i];
    }
    delete[] dop;
    delete[] dopSizes;
}

void printArr(int arr[], int n) {
    for (int i = 0; i < n; i++) {
        cout << arr[i] << " ";
    }
    cout << endl;
}

int findMaxRazryad(int arr[], int n) {
    int maxVal = arr[0];
    for (int i = 1; i < n; i++) {
        if (arr[i] > maxVal) maxVal = arr[i];
    }
    int razryad = 0;
    while (maxVal > 0) {
        maxVal /= 10;
        razryad++;
    }
    return razryad;
}

int findk(int arr[], int n) {
    int maxVal = arr[0];
    for (int i = 1; i < n; i++) {
        if (arr[i] > maxVal) {
            maxVal = arr[i];
        }
    }
    return maxVal + 1;
}

int main() {
    setlocale(LC_ALL, "RUS");
    cout << "сортировка подсчетом" << endl;
    int arr[] = { 4, 2, 2, 8, 3, 3, 1, 0, 4, 0, 5, 7 };
    int n1 = sizeof(arr) / sizeof(arr[0]);
    int k1 = findk(arr, n1);

    cout << "исходный массив: ";
    printArr(arr, n1);

    countingSort(arr, n1, k1);

    cout << "результат countingSort: ";
    printArr(arr, n1);
    cout << endl;

    cout << "поразрядная сортировка LSD" << endl;
    int arr2[] = { 160, 60, 144, 12, 78, 56, 2, 122 };
    int n2 = sizeof(arr2) / sizeof(arr2[0]);
    int k2 = 10;
    int m = findMaxRazryad(arr2, n2);

    cout << "исходный массив: ";
    printArr(arr2, n2);

    radixSortLSD(arr2, n2, m, k2);

    cout << "результат radixSortLSD: ";
    printArr(arr2, n2);

    return 0;
}
\end{Verbatim}

\end{document}
